% Question 13
\question
In a workshop, $2$ identical solid metal right circular cylinders of base radius $R\text{ cm}$ are melted and recast into $27$ smaller identical solid right circular cylinders of base radius $r\text{ cm}$ and height $10\text{ cm}$. It is given that the base area of a larger circular cylinder is $9$ times that of a smaller one.
\begin{parts}
\part Find
\begin{subparts}
\subpart $r : R$,
\subpart the height of a larger circular cylinder. \\*
\null\rmarks{5}
\end{subparts}
\part A craftsman claims that a smaller circular cylinder and a larger circular cylinder are similar. Do you agree? Explain your answer. \rmarks{2}
\end{parts}

\begin{solution}
\begin{parts}
\part
\begin{subparts}
\subpart 
\begin{align*}
\frac{\pi r^2}{\pi R^2} &= \frac{1}{9} \quad &\text{(1M)} \\
\left(\frac{r}{R}\right)^2 &= \frac{1}{9} \\
\frac{r}{R} &= \frac{1}{3} \\
r : R &= 1 : 3 \quad &\text{(1A)}
\end{align*}
\subpart Let $h\text{ cm}$ be the height of a larger circular cylinder.
\begin{align*}
2\pi R^2 h &= 27(\pi r^2(10)) \quad &\text{(1M)} \\
h &= \frac{270}{2}\left(\frac{r}{R}\right)^2 \\
h &= \frac{270}{2}\left(\frac{1}{3}\right)^2 \quad &\text{(1M)} \\
h &= 15 \quad &\text{(1A)}
\end{align*}
Thus, the height of a larger circular cylinder is $15\text{ cm}$.
\end{subparts}
\part Note that
\[
\frac{\text{The base radius of a smaller circular cylinder}}{\text{The base radius of a larger circular cylinder}} = \frac{1}{3} \quad \text{and} \quad \frac{\text{The height of a smaller circular cylinder}}{\text{The height of a larger circular cylinder}} = \frac{10}{15} = \frac{2}{3}
\]
Since
\[
\frac{\text{The base radius of a smaller circular cylinder}}{\text{The base radius of a larger circular cylinder}} \ne \frac{\text{The height of a smaller circular cylinder}}{\text{The height of a larger circular cylinder}} \quad \text{(1M)}
\]
the two circular cylinders are not similar. \\
Thus, the claim is disagreed. \hfill (1A f.t.)
\end{parts}
\end{solution}